\documentclass[11pt]{beamer}

\usetheme{Madrid}
\usecolortheme{default}
\definecolor{unicaucaverde}{RGB}{0,90,60}
\setbeamercolor{structure}{fg=unicaucaverde}
\setbeamertemplate{navigation symbols}{}

\usepackage[utf8]{inputenc}
\usepackage[spanish]{babel}
\usepackage{amsmath,amssymb}
\usepackage{listings}
\usepackage{xcolor}

\lstdefinestyle{cstyle}{
  language=C,
  basicstyle=\ttfamily\scriptsize,
  keywordstyle=\color{unicaucaverde}\bfseries,
  commentstyle=\color{gray}\itshape,
  stringstyle=\color{orange!70!black},
  breaklines=true,
  showstringspaces=false,
  tabsize=2
}
\lstset{style=cstyle}

\title[Semana 4: Ámbito avanzado y recursión]{Ámbito avanzado y recursión}
\subtitle{Programación --- Semana 4 de 16}
\author{Universidad del Cauca}
\date{2026}

\begin{document}

\begin{frame}
  \titlepage
\end{frame}

\begin{frame}{Semana 4: dónde estamos}
  \begin{itemize}
    \item Continuación de Funciones y modularidad (Semana 3): RA3, parte 2 de 2.
    \item Subtemas 4.1--4.2 --- RA3.3, RA3.4.
    \item 4 horas: variables estáticas locales, luego recursión y pila de llamadas.
  \end{itemize}
  \tableofcontents
\end{frame}

\section{Variables estáticas y ámbito avanzado}

\begin{frame}{Ámbito frente a tiempo de vida}
  \begin{itemize}
    \item \textbf{Ámbito}: dónde es visible el nombre de la variable.
    \item \textbf{Tiempo de vida}: cuánto tiempo existe la variable en memoria.
    \item En una variable local ordinaria coinciden; \texttt{static} los separa.
  \end{itemize}
  \begin{block}{Variable estática local}
    Ámbito: solo dentro de su función. Tiempo de vida: todo el programa. Se inicializa una única
    vez y conserva su valor entre llamadas.
  \end{block}
\end{frame}

\begin{frame}[fragile]{\texttt{static}: contador que persiste}
\begin{lstlisting}
int contadorLlamadas(void) {
    static int veces = 0;   // se inicializa UNA sola vez
    veces++;
    return veces;
}

int main(void) {
    printf("%d\n", contadorLlamadas());   // 1
    printf("%d\n", contadorLlamadas());   // 2
    printf("%d\n", contadorLlamadas());   // 3
    return 0;
}
\end{lstlisting}
  \textbf{Taller (hora 2):} \texttt{contadorLlamadas}, \texttt{siguienteId}.
\end{frame}

\begin{frame}[fragile]{Bloques anidados y sombreado}
\begin{lstlisting}
int main(void) {
    int x = 1;
    {
        int x = 2;   // sombrea (oculta) la x externa
        printf("%d\n", x);   // 2
    }
    printf("%d\n", x);       // 1: la x externa no cambio
    return 0;
}
\end{lstlisting}
  \begin{alertblock}{Cuidado}
    Combinar \texttt{static} con nombres repetidos entre bloques o funciones es fuente de errores
    difíciles de rastrear.
  \end{alertblock}
\end{frame}

\section{Recursión}

\begin{frame}[fragile]{Caso base y caso recursivo}
\begin{lstlisting}
int factorial(int n) {
    if (n <= 1) {                  // caso base
        return 1;
    }
    return n * factorial(n - 1);   // caso recursivo
}
\end{lstlisting}
  \begin{itemize}
    \item \textbf{Caso base}: se resuelve directamente, sin más llamadas.
    \item \textbf{Caso recursivo}: reduce el problema y se apoya en la llamada más pequeña.
    \item Sin caso base alcanzable $\Rightarrow$ recursión infinita $\Rightarrow$
          \emph{stack overflow}.
  \end{itemize}
\end{frame}

\begin{frame}{La pila de llamadas}
  \begin{itemize}
    \item Cada llamada activa apila un \textbf{marco} (parámetros + variables locales).
    \item El marco se retira al terminar y devolver su resultado.
    \item En \texttt{factorial(5)}: los marcos de \texttt{factorial(5..1)} coexisten mientras las
          llamadas están pendientes.
  \end{itemize}
\end{frame}

\begin{frame}[fragile]{Recursión frente a iteración: Fibonacci}
\begin{lstlisting}
int fibonacciRec(int n) {
    if (n == 0) return 0;
    if (n == 1) return 1;
    return fibonacciRec(n-1) + fibonacciRec(n-2);
}

int fibonacciIter(int n) {
    if (n == 0) return 0;
    int a = 0, b = 1;
    for (int i = 2; i <= n; i++) { int s = a+b; a = b; b = s; }
    return b;
}
\end{lstlisting}
  \begin{alertblock}{Clave}
    Recursiva: más clara, refleja la definición matemática. Iterativa: más eficiente en tiempo y
    memoria. La elección es de diseño, no una regla fija.
  \end{alertblock}
  \vfill
  \textbf{Taller (hora 4):} \texttt{potencia}, \texttt{sumaHasta}, \texttt{fibonacci} recursivas
  y sus versiones iterativas.
\end{frame}

\begin{frame}{Cierre de la semana}
  \begin{itemize}
    \item RA3.3: variables estáticas locales, ámbito frente a tiempo de vida.
    \item RA3.4: recursión, caso base/recursivo, pila de llamadas, recursión frente a iteración.
    \item Taller calificado al cierre de la Hora 4.
  \end{itemize}
\end{frame}

\end{document}
